\subsection{Two-run forking and quantitative knowledge error}\label{sec:mv:forking}
For soundness the \emph{honest} sigma verifier
chooses the bit challenges uniformly.
This is a different security game from the
malicious-verifier ZK simulation above.

\begin{definition}[Resettable fixed-statement prover]\label{def:mv:reset-prover}
Fix a public OR statement $(C,S)$ before
the first one-bit verifier challenge.
A resettable prover has a fixed initial
state and a complete fixed random tape
for a sequence of $\kappa$ fresh one-bit
mini-proofs for that same statement.
Given any verifier challenge string
$\mathbf e=(e_1,\ldots,e_\kappa)
\in\{0,1\}^\kappa$,
the resulting messages and accepting
or rejecting outcome are deterministic.
The prover may adapt its next first
message to earlier challenges, but not
to the current challenge before sending
that first message.
A reset restarts it from the same
initial state and random tape.
\end{definition}

\begin{definition}[Accepting challenge set]\label{def:mv:accepting-set}
For the fixed deterministic state of
Definition~\ref{def:mv:reset-prover},
let
\[
 \mathcal A:=\{\mathbf e\in\{0,1\}^{\kappa}:
       \text{every mini-proof accepts}\},\quad
 m:=|\mathcal A|,\quad p:=m/2^\kappa.
\]
The quantity $p$ is the conditional success
probability against independent uniform
verifier challenge bits.
\end{definition}

\begin{lemma}[Earliest difference yields one OR fork]\label{lem:mv:earliest}
If two distinct challenge strings
$\mathbf e,\mathbf e'\in\mathcal A$
are both accepted from the same
reset prover state and random tape,
then their transcripts contain two
accepting responses to the
\emph{same} OR first message
under the \emph{different} global
bits $0$ and $1$.
\end{lemma}
\begin{proof}
Let $i$ be the smallest index
with $e_i\ne e_i'$.
All earlier verifier challenges
and prover responses are identical,
because the starting state and
random tape are fixed and the
prover is causal.
Immediately before the $i$-th
challenge, both runs therefore
have the same prover state
and the same first message $T_i$.
Their current challenges differ:
one is $0$ and the other $1$.
Both full challenge strings lie in
$\mathcal A$, so the verifier
accepts the $i$-th mini-proof
in both runs. The two local
accepting transcripts have the
required common first message
and distinct global challenges.
\end{proof}

\begin{theorem}[Combinatorial knowledge extraction from two accepting runs]\label{thm:mv:two-run-extract}
Under the hypotheses of
Lemma~\ref{lem:mv:earliest},
one can compute in polynomial
time from the two transcripts
an OR witness $(t_j,r)$ such that
\[
 \boxed{t_j\in S,\qquad C=g^{t_j}h^r.}
\]
\end{theorem}
\begin{proof}
The preceding lemma supplies two
accepting mini-proofs with the same
statement and first message and
two different global challenges.
Apply the explicit extraction
formula of
Theorem~\ref{thm:realize:or-extract}.
It produces a branch index $j$
and its opening exponent $r$.
Comparing the two transcripts to
find the earliest challenge
difference, and computing the
one modular inverse, are
polynomial-time operations.
\end{proof}

\begin{lemma}[Exact two-run forking probability]\label{lem:mv:fork-prob}
For a fixed deterministic prover
with success probability
$p=m/2^\kappa$, sample two independent
uniform challenge strings and reset
the same prover state between runs.
The probability that both runs
accept and their challenge strings
are distinct is exactly
\[
 \boxed{p(p-2^{-\kappa}).}
\]
\end{lemma}
\begin{proof}
There are $m$ accepting strings.
The probability that both independent
samples belong to $\mathcal A$ is
$(m/2^\kappa)^2=p^2$.
Among those ordered pairs,
there are exactly $m$ diagonal
pairs with identical challenge strings.
Each ordered pair has probability
$2^{-2\kappa}$, so the diagonal
probability equals
$m2^{-2\kappa}=p2^{-\kappa}$.
Subtracting yields the formula.
\end{proof}

\begin{theorem}[Randomized-prover two-run forking bound]\label{thm:mv:fork-average}
Let a randomized causal prover have
acceptance probability $\varepsilon$
against $\kappa$ independent uniform
challenge bits, on a fixed statement.
Run it twice from the \emph{same}
initial state and \emph{same} private
random tape but with independent
verifier challenge strings.
The probability of obtaining an
extractable fork is at least
\[
 \boxed{
 \max\{0,\varepsilon(\varepsilon-2^{-\kappa})\}.}
\]
\end{theorem}
\begin{proof}
Condition on the prover's
initial private random tape $R$,
and write $p_R$ for its conditional
acceptance probability.
Lemma~\ref{lem:mv:fork-prob}
gives conditional fork probability
$p_R^2-2^{-\kappa}p_R$.
Averaging over $R$ yields
\[
 \mathbb E[p_R^2]
 -2^{-\kappa}\mathbb E[p_R].
\]
By the elementary nonnegativity
of variance,
$\mathbb E[p_R^2]\ge
(\mathbb E[p_R])^2=\varepsilon^2$.
The expectation of $p_R$ is
$\varepsilon$ by the law of total
probability. The lower bound
is therefore
$\varepsilon^2-2^{-\kappa}\varepsilon$.
A probability is also nonnegative,
giving the stated maximum.
Theorem~\ref{thm:mv:two-run-extract}
turns any successful fork into
a valid OR witness.
\end{proof}

\begin{corollary}[Knowledge-error dichotomy]\label{cor:mv:knowledge-error}
Fix a resettable prover running in
polynomial time. Either
$\varepsilon\le2^{-\kappa}+\delta$,
or a two-run black-box extractor
finds an OR witness with probability
at least $\varepsilon\delta$,
where $\delta>0$ is any chosen
threshold.
\end{corollary}
\begin{proof}
If the first alternative fails,
then $\varepsilon-2^{-\kappa}>\delta$.
Theorem~\ref{thm:mv:fork-average}
makes the extraction probability
at least
$\varepsilon(\varepsilon-2^{-\kappa})
>\varepsilon\delta$.
\end{proof}

\begin{lemma}[Why a single global $\kappa$-bit guess is costly]\label{lem:mv:parallel-cost}
Suppose a verifier independently
samples a uniform bit vector
$\mathbf e\in\{0,1\}^{\kappa}$
\emph{after} receiving all first
messages, and a simulator guesses
one independent uniform
$\widehat{\mathbf e}$ in advance.
Then the probability of a complete
match is exactly $2^{-\kappa}$,
so the expected number of complete
independent trials until a match
is $2^\kappa$.
\end{lemma}
\begin{proof}
For every guessed vector, precisely
one of the $2^\kappa$ equally likely
verifier vectors matches.
Independent repeated trials
therefore have success probability
$2^{-\kappa}$ and geometric
expected trial count its reciprocal,
namely $2^\kappa$.
\end{proof}

\begin{remark}[What this theorem actually extracts]
Theorems~\ref{thm:mv:two-run-extract}
and~\ref{thm:mv:fork-average}
establish extraction for \emph{one fixed
OR statement} across resettable
one-bit repetitions. They do not
by themselves yield a single
online extractor for every
interleaved proof and a complete
Ring Switching transcript.
That multistage composition must
respect the two outer challenge
boundaries; we state it separately.
\end{remark}
